\section{RNN 的优势与局限}

\subsection{优势}

\begin{figure}[H]
\centering
\includegraphics[width=0.95\textwidth]{slides-images/slide-034.jpg}
\caption{RNN 的核心优势\protect\footnotemark}
\end{figure}
\footnotetext{Slides 第35页。}

\begin{enumerate}
    \item \textbf{可处理任意长度的输入}：没有上下文长度限制（与 Transformer 形成对比）
    \item \textbf{理论上可利用任意远的历史信息}：如果隐藏状态有效地编码了历史，计算可以依赖很久以前的输入
    \item \textbf{模型大小不随输入长度增长}：参数数量固定，只由隐藏状态维度决定
    \item \textbf{每个时间步应用相同的变换}：对称性和概念简洁性
\end{enumerate}

\subsection{局限}

\begin{enumerate}
    \item \textbf{训练时难以并行化}：计算 $h_t$ 必须先计算 $h_{t-1}$，这是一个严格的顺序依赖
    \item \textbf{实际上难以访问远距离信息}：固定大小的隐藏状态无法无损地编码任意长的历史
    \item \textbf{梯度消失/爆炸}：长序列的反向传播中，梯度会指数级衰减或增长
\end{enumerate}

\begin{warningbox}{训练的不可并行性是 RNN 的致命弱点}
现代深度学习的成功很大程度上依赖于 GPU 的大规模并行计算能力。RNN 的顺序依赖结构使得它无法充分利用 GPU 的并行性——这正是 Transformer 最终取代 RNN 的核心原因之一。在训练时，Transformer 可以同时处理序列中的所有位置，而 RNN 必须逐步计算。
\end{warningbox}

\subsection{本章小结}

RNN 在理论上具有处理无限长序列的能力，但在实践中受到训练效率（不可并行）和梯度流动（消失/爆炸）的严重限制。这些局限性推动了 LSTM、GRU 以及最终 Transformer 的诞生。

%% ========== 计算机视觉应用 ==========

